% Sections 1.1-1.6, from lectures/L01/appendix/a_combinational_logic.md (A.1-A.6).

\section{Digital signals and logic gates}
\label{c1:sec:gates}

A digital signal takes one of two values, \code{0} or \code{1}, instead of the continuous range an
analogue signal uses. Anything close enough to \code{0} or \code{1} reads as exactly \code{0} or
\code{1}, so small amounts of electrical noise change nothing. That noise immunity is the reason
nearly all modern computation, from microcontrollers to FPGAs, is digital.

A \textbf{logic gate} has one or more single-bit inputs, a single-bit output and a fixed Boolean
function between them.

\textbf{Combinational logic}, this chapter's subject, produces outputs from the current inputs
alone, with no memory. \textbf{Sequential networks} (\chapref{c3:ch}) depend on the current inputs
\emph{and} on state stored in flip-flops.

Every basic two-input gate, plus the unary \code{NOT}:

\begin{narrowtable}{cc*{6}{c}}
  \thead{A} & \thead{B} & \thead{AND} & \thead{OR} & \thead{NAND} & \thead{NOR} & \thead{XOR} &
    \thead{XNOR} \\
  \midrule
  0 & 0 & 0 & 0 & 1 & 1 & 0 & 1 \\
  0 & 1 & 0 & 1 & 1 & 0 & 1 & 0 \\
  1 & 0 & 0 & 1 & 1 & 0 & 1 & 0 \\
  1 & 1 & 1 & 1 & 0 & 0 & 0 & 1 \\
\end{narrowtable}

\begin{narrowtable}{cc}
  \thead{A} & \thead{NOT A} \\
  \midrule
  0 & 1 \\
  1 & 0 \\
\end{narrowtable}

Worth committing to memory straight from the tables:
\begin{itemize}
  \item \code{NAND}, \code{NOR} and \code{XNOR} are the inverses of \code{AND}, \code{OR} and
    \code{XOR}.
  \item \code{XOR} is \code{1} when the inputs differ. Beyond two inputs it generalises to
    \textbf{odd parity}: \code{1} when an odd number of inputs are \code{1}
    (Exercise~\ref{c1:ex:xor3}).
  \item Every Boolean function can be built from \code{AND}, \code{OR} and \code{NOT}, and also
    from \code{NAND} alone or \code{NOR} alone.
\end{itemize}

% ------------------------------------------------------------------------------------------
\section{Boolean algebra as notation}
\label{c1:sec:boolean}

Three operators, which map straight onto the gates above:

\begin{booktable}{@{}p{5em}L@{}}
  \thead{Gate} & \thead{Notation in Boolean algebra} \\
  \midrule
  AND & multiplication: $A$ \textbf{AND} $B$ is written $AB$ \\
  OR & addition: $A$ \textbf{OR} $B$ is written $A + B$ \\
  NOT & prime: \textbf{NOT} $A$ is written $A'$ \\
\end{booktable}

A function written as an OR of AND terms, for example $X = AB' + CD$, is in \textbf{sum-of-products}
form.

Every truth table can be read off directly as a sum of products:
\begin{enumerate}
  \item Find every row where the output is \code{1}.
  \item For each of them, write an AND term containing every input: plain if the input is \code{1}
    on that row, primed if it is \code{0}.
  \item OR those terms together.
\end{enumerate}

For example:

\begin{narrowtable}{ccc}
  \thead{A} & \thead{B} & \thead{X} \\
  \midrule
  0 & 0 & 0 \\
  0 & 1 & 1 \\
  1 & 0 & 1 \\
  1 & 1 & 0 \\
\end{narrowtable}

The row $AB = 01$ contributes $A'B$ and the row $AB = 10$ contributes $AB'$, which gives
$X = A'B + AB'$. Check it against the table: that is \code{XOR}, so the same network is a single
\code{XOR} gate instead of two \code{AND} gates, two \code{NOT} gates and an \code{OR}.

\subsection*{The laws, as a reminder}
You have met them before. They are here so that the book has a single notation for them and so
that later chapters can refer to them, not because they need teaching:

\begin{booktable}{@{}p{8em}p{9em}L@{}}
  \thead{Law} & \thead{With OR} & \thead{With AND} \\
  \midrule
  Identity & $A + 0 = A$ & $A \cdot 1 = A$ \\
  Dominance & $A + 1 = 1$ & $A \cdot 0 = 0$ \\
  Idempotence & $A + A = A$ & $A \cdot A = A$ \\
  Complement & $A + A' = 1$ & $A \cdot A' = 0$ \\
  Absorption & $A + AB = A$ & $A(A + B) = A$ \\
  De Morgan & $(A + B)' = A'B'$ & $(AB)' = A' + B'$ \\
\end{booktable}

\textbf{De Morgan is the law this book leans on}, less for minimisation than as a rule for how an
inversion moves through a gate: an \code{OR} with both inputs inverted \emph{is} a \code{NAND}, and
an \code{AND} with both inputs inverted \emph{is} a \code{NOR}. That is why a \code{NAND} gate can
build anything, which Exercise~\ref{c1:ex:nand} lets you write in VHDL, and it comes back every
time an active-low signal meets logic written for active-high. In a course where every reset and
every button is active-low, that is often.

The gap between a directly read sum of products and the smallest network is what a \textbf{Karnaugh
map} closes, and \secref{c2:sec:kmap} works through one from start to finish. Not because the
technique is new to you, but because it is the minimised equation that gets written in VHDL, and
\chapref{c8:ch} derives its state machine's next-state logic in exactly that way.

% ------------------------------------------------------------------------------------------
\section{Building and simulating a circuit in CircuitVerse}
\label{c1:sec:cv}

CircuitVerse (\url{https://circuitverse.org/simulator}) is a free, browser-based logic simulator.
Every gate network in this book is built and tested there by hand before it is written in VHDL:
\begin{enumerate}
  \item Open the simulator and start a new project.
  \item Place one \textbf{input} element per input signal and one \textbf{output} element per
    output signal, renamed to match your variables.
  \item Place the gates your equation calls for, from the \emph{Gates} panel.
  \item Wire the network up one gate at a time, following your derived equation. You should
    already know exactly which gates you need before you open the simulator, which is why
    \secref{c1:sec:boolean} comes first.
  \item Toggle every input and check that the output matches your truth table on \textbf{every}
    row.
  \item Save with \code{Project -> Save Online}, or \code{Project -> Download as image/data} for a
    re-importable \file{.cv} file.
\end{enumerate}

\textbf{Step 5 is not optional, and nothing else does it for you.} Three rules make it a check
rather than a formality:
\begin{itemize}
  \item \textbf{Exhaustive, not sampled.} Two inputs are four rows, three are eight. At those
    sizes there is no excuse for sampling. For the sequential networks further on this becomes
    ``step the clock and check at every edge'', the same discipline applied to time.
  \item \textbf{Derive first, simulate second.} Draw the circuit and \emph{then} decide what it
    should do, and you will read your own expectations out of the simulation, which proves
    nothing.
  \item \textbf{A discrepancy is information.} It says the drawing and the equation disagree, and
    that one of them is wrong.
\end{itemize}

Every hardware bug you will chase is found by predicting a behaviour and then checking it. Doing it
by hand on a circuit with four gates is how you learn to do it on one you cannot take in at a
glance.

Once the design is VHDL, a self-checking testbench takes that job over. Every worked example and
every exercise has one handed out with it; you run them rather than write them, which
\secref{c2:sec:testbenches} works through.

% ------------------------------------------------------------------------------------------
\section{Exactly enough VHDL that you can read and write a gate network}
\label{c1:sec:vhdl}

VHDL describes hardware, not a sequence of instructions. Every statement below runs concurrently
and continuously, exactly like the gates it describes. The language arrives as it is needed; this
section is enough for a simple combinational network.

Every VHDL module has two parts:
\begin{itemize}
  \item An \textbf{entity}, the view from outside: the ports, and nothing about what happens
    inside.
  \item An \textbf{architecture}, the implementation behind it.
\end{itemize}

That split is a contract, not ceremony. The entity is everything another design needs in order to
use the module, which is what lets one design instantiate another without reading it, and what lets
you rewrite an implementation without touching what depends on it. The shape is a header beside a
source file, or an interface beside its class; the difference is that VHDL requires it for every
module, with no way to skip the interface or to leak the implementation through it.

Take \code{or_gate}: the inputs \code{a}, \code{b}, the output \code{x}, and $x = a + b$.

\begin{vhdlcode}
entity or_gate is
    port(a, b: in std_logic;
         x   : out std_logic);
end entity;
\end{vhdlcode}

\lecfig[0.62]{lectures/L01/appendix/images/or_gate_entity.png}{The entity \code{or_gate}: the view
from outside, and nothing about what happens inside.}{c1:fig:entity}

\code{std_logic} is the signal type for a single bit. It is not built into VHDL but comes from the
package \code{std_logic_1164}, which is why every file in this book opens with:

\begin{vhdlcode}
library ieee;
use ieee.std_logic_1164.all;
\end{vhdlcode}

It has nine values instead of two, because real hardware needs more than true and false:
\begin{itemize}
  \item \code{'0'} and \code{'1'}: driven low and driven high. The only two you will write in this
    book.
  \item \code{'Z'}: high impedance, nothing is driving the wire at all.
  \item \code{'U'}: uninitialised, what a signal holds before anything has assigned it a value.
  \item \code{'X'}: unknown, what a simulator shows when two things drive the same signal and
    disagree.
  \item \code{'-'}, \code{'W'}, \code{'L'}, \code{'H'}: don't-care values and weakly driving
    values, rarely written by hand.
\end{itemize}

Recognising the seven you will not write still matters: they are how a simulator tells you
something is wrong. A signal left sitting at \code{'U'} or \code{'X'} is a bug, not a value.

The architecture fills the box in:

\begin{vhdlcode}
architecture behaviour of or_gate is
begin
    x <= a or b;
end architecture;
\end{vhdlcode}

\lecfig[0.62]{lectures/L01/appendix/images/or_gate_arch.png}{The architecture \code{or_gate}: the
implementation behind the entity.}{c1:fig:arch}

Together they make up the complete module:

\lecfig[0.62]{lectures/L01/appendix/images/or_gate_module.png}{The module \code{or_gate}: entity and
architecture together.}{c1:fig:module}

\code{x <= a or b;} is \secref{c1:sec:boolean}'s Boolean algebra almost verbatim: VHDL's
\code{and}, \code{or}, \code{not} and \code{xor} work directly on \code{std_logic}, so every
equation you derived by hand becomes VHDL with hardly any translation at all.

The assignment is \textbf{concurrent}. It does not run once; it holds continuously, like a wire
between real gates.

The keyword \code{signal} declares an internal wire inside an architecture, as opposed to a port
visible from outside. You use one in Exercise~\ref{c1:ex:adas} to name an intermediate result
instead of nesting it in place, and \chapref{c2:ch} leans on the same idea throughout.

% ------------------------------------------------------------------------------------------
\section{The complete module: \code{or_gate}}
\label{c1:sec:module}

\Secref{c1:sec:vhdl}'s parts together make a complete, synthesisable file. This is the whole of
\file{or_gate.vhd}, and it is a real design that synthesises and runs on the DE0-CV:

\vhdlfile{lectures/L01/or_gate/or_gate.vhd}

\begin{itemize}
  \item The \code{library}/\code{use} clauses pull in \code{std_logic}.
  \item \code{entity} declares the outside: two inputs, one output.
  \item \code{architecture} implements the inside, as a single concurrent assignment.
\end{itemize}

That is the smallest complete VHDL design there is, and the shape never changes. Every module in
this book, all the way up to the CAN controller in \chapref{c17:ch}, is the same
\code{library}/\code{entity}/\code{architecture} skeleton. What grows is the architecture body,
never the frame.

\paragraph{Checking your work}
\code{or_gate} comes with an \file{or_gate_tb.vhd}, which drives all four input combinations and
checks the output:

\begin{shell}
cd lectures/L01/or_gate
ghdl -a --std=93 or_gate.vhd or_gate_tb.vhd
ghdl -e --std=93 or_gate_tb
ghdl -r --std=93 or_gate_tb --assert-level=error --stop-time=10ms
\end{shell}

No output beyond a closing note means every check passed. Testbenches are treated properly in
\secref{c2:sec:testbenches}; until then, take these three commands as the way you confirm that a
module works.

% ------------------------------------------------------------------------------------------
\section{A brake assistant, from requirement to circuit}
\label{c1:sec:adas}

The session builds a different circuit from \code{or_gate}, live, so that you get to see the same
path walked twice: once here in writing, and once on something you have not already read the answer
to.

The requirement: \textbf{the car brakes if the driver brakes, or if the assistance system decides
to.} The assistance system decides to when either detector sees something, unless it is at the same
time reporting a fault, in which case it is not trusted to brake at all.

\begin{booktable}{@{}p{8em}p{4em}L@{}}
  \thead{Port} & \thead{Direction} & \thead{Meaning} \\
  \midrule
  \code{driver_brake} & in & The driver is on the brake pedal. \\
  \code{sensor} & in & The distance sensor sees an obstacle. \\
  \code{radar} & in & The radar sees an obstacle. \\
  \code{error} & in & The assistance system reports a fault. \\
  \code{engine_brake} & out & Apply the brakes. \\
\end{booktable}

Three sentences together determine the whole structure:
\begin{itemize}
  \item either detector is enough on its own, and neither of them outweighs the other.
  \item a fault raises no alarm; it \emph{prevents} the assistance system from being able to brake.
  \item the driver can always brake, even when the assistance system is broken and doing nothing.
\end{itemize}

The last is a safety requirement, and it is the one to hold on to as you derive the network: the
pedal is not an input to the assistance logic, it goes past it. Work out what changes if the pedal
is instead fed through the fault logic, and you have found the bug the structure exists to prevent.

\subsection*{The truth table}
Sixteen rows, in the order a four-bit counter would produce them. \code{adas_brake} is not a port;
it is the assistance system's own decision, shown as an intermediate column so that you can see
where the structure comes from.

\begin{narrowtable}{cccccc}
  \thead{\code{driver_brake}} & \thead{\code{sensor}} & \thead{\code{radar}} &
    \thead{\code{error}} & \thead{\code{adas_brake}} & \thead{\code{engine_brake}} \\
  \midrule
  0 & 0 & 0 & 0 & 0 & 0 \\
  0 & 0 & 0 & 1 & 0 & 0 \\
  0 & 0 & 1 & 0 & 1 & 1 \\
  0 & 0 & 1 & 1 & 0 & 0 \\
  0 & 1 & 0 & 0 & 1 & 1 \\
  0 & 1 & 0 & 1 & 0 & 0 \\
  0 & 1 & 1 & 0 & 1 & 1 \\
  0 & 1 & 1 & 1 & 0 & 0 \\
  1 & 0 & 0 & 0 & 0 & 1 \\
  1 & 0 & 0 & 1 & 0 & 1 \\
  1 & 0 & 1 & 0 & 1 & 1 \\
  1 & 0 & 1 & 1 & 0 & 1 \\
  1 & 1 & 0 & 0 & 1 & 1 \\
  1 & 1 & 0 & 1 & 0 & 1 \\
  1 & 1 & 1 & 0 & 1 & 1 \\
  1 & 1 & 1 & 1 & 0 & 1 \\
\end{narrowtable}

The lower half is \code{1} throughout: as soon as \code{driver_brake} is set, nothing else changes
the answer. That is what ``overrides everything'' looks like in a table.

\subsection*{Four answers to bring to the session}
Everything above is the specification. The derivation, the gate network and the VHDL code are what
gets built live, and it is worth arriving there with an answer of your own to compare against
rather than one you have already read. So, from the table above:
\begin{itemize}
  \item Read a sum-of-products equation for \code{engine_brake} straight out of it, the same way
    \secref{c1:sec:boolean} did.
  \item Simplify it by hand. The requirement is three sentences long, so the minimal network is
    considerably smaller than what the row-by-row reading gives you. Work out how much smaller
    before you are told.
  \item Count the gates your simplified equation needs, and write the number down.
  \item Sketch the network, and decide which single gate the safety argument above rests on.
\end{itemize}

Bring those four answers to the session. Being wrong about one of them is worth more than not
having guessed, because you then know exactly which step to look at again when the live derivation
reaches it. It is \secref{c1:sec:cv}'s ``derive first, simulate second'' applied to a circuit you
have not built yet, and it is the habit the rest of the book is built on. Afterwards you write the
VHDL code yourself, in Exercise~\ref{c1:ex:adas}, and a testbench decides whether your equations and
the truth table above agree.

\subsection*{Why there is a testbench for it}
Sixteen combinations are few enough to check by hand once. They are not few enough to check every
time someone edits the file, which is the situation every module in this book is in from here on.

A \textbf{testbench} is a second VHDL file that does the checking. It is not part of the design and
never reaches the FPGA: it drives the inputs through every case, works out what the output should
have been, compares and reports. For \code{adas}, which is written alongside it during the session,
that is a loop over the sixteen combinations restating the braking rule independently.

\emph{Independently} is the important word. A testbench that read its expected answer out of the
design would agree with it by construction, even when the design is wrong. Restating the
requirement separately is what makes it a check.

That is why nearly every directory in the course repository has a \file{_tb.vhd} beside the module:
it is how you can be handed a module, change it, and know within a second whether you have broken
it. \Secref{c2:sec:testbenches} works through how they are run; you are not asked to write one
anywhere in this book.
